\documentclass[10pt,a4paper]{amsart}
\usepackage[margin=19mm]{geometry}
\usepackage{amsmath,amssymb,mathtools}
\usepackage[scr=rsfs,cal=euler]{mathalfa}
\usepackage{xfrac}
\usepackage[unicode,hidelinks,bookmarks=false]{hyperref}
\newcommand{\ie}{i.e.\ }
\newcommand{\cf}{cf.\ }
\newcommand{\resp}{respectively\ }
\newcommand{\Subsection}{\subsection}
\providecommand{\nobreakdash}{\nobreak\hbox{-}}
\setlength{\emergencystretch}{2em}
\multlinegap0pt
\usepackage[shortlabels]{enumitem}
\usepackage{tensor}
\usepackage[all]{xy}
\usepackage{xcolor}
\usepackage{tikz}
\usepackage{mathtools}
\newtheorem{lemma}{Lemma}
\newcommand{\PP}{\mathbb{P}}
\newcommand{\Pic}{\operatorname{Pic}}
\newcommand{\Q}{\mathbb{Q}}
\newcommand{\Z}{\mathbb{Z}}
\title{On Proper Descent of Smooth Affine Surfaces\\with Finite Homotopy Rank-Sum}
\author{Buddhadev Hajra}
\date{Source: arXiv:2601.17479v1 (CC BY 4.0)}
\begin{document}
\maketitle

\noindent\emph{Source and licence.} On Proper Descent of Smooth Affine Surfaces\\with Finite Homotopy Rank-Sum. Version arXiv:2601.17479v1 (CC BY 4.0), distributed by arXiv under the Creative Commons Attribution 4.0 International licence, http://creativecommons.org/licenses/by/4.0/, as recorded in the arXiv licence metadata for this exact version and shown on the abstract page https://arxiv.org/abs/2601.17479v1. Full licence notice: \texttt{licenses/arxiv-2601.17479-CC-BY-4.0.txt}.\par\smallskip

\noindent\textit{Editorial scope.} Counted unit: the article's lemma fibration (source0.tex lines 496--502). The article's complete original proof of it is delivered with it (source0.tex lines 504--522, one \texttt{proof} environment of 19 lines). No mathematics is written, repaired or re-derived: the statement and the proof are the article's own lines, in the article's own notation, and no retained line is substituted or re-flowed.  No further source-local context is needed: every cross-reference of every retained line resolves inside the two delivered spans.  Nothing was omitted from any retained span, so the package carries no omission record.  No retained line contains a citation, so the wrapper carries no bibliography and no bibliography entry.  No retained line contains a figure environment, an image-inclusion command or drawing code, so no image is delivered and the package declares no asset dependency.  Editorial formatting only, in addition: an AMS \texttt{amsart} preamble that restates the article's own declarations and macros used by the retained lines, namely the environments \texttt{lemma} on separate counters, together with the standard packages those lines need.  The retained environments print in this document as Lemma 1 in order of appearance, whereas the article prints the same items with the numbering of its own full document.  The complete native source of the article as distributed in the arXiv e-print for this exact version is bundled unchanged as \texttt{sources/193316/source0.tex}.
\begin{lemma}\label{Lem: Q-factorial fibration}
Let $f: X \to B$ be an $F$-fibration from a smooth affine surface $X$ onto a smooth algebraic curve $B$, where $F$ is a smooth affine rational curve. If $X$ is $\Q$-factorial, then the following hold:
\begin{enumerate}[\indent\rm(1)]
    \item $B$ is affine.
    \item All fibers of $f$ are irreducible. 
\end{enumerate}
\end{lemma}

\begin{proof}
Let $f : X \to B$ be an $F$-fibration as in the statement, with $X$ a smooth affine $\Q$-factorial surface and $B$ a smooth algebraic curve. Embed $X$ as an open subset of a smooth projective surface $V$ such that $f$ extends to a $\PP^1$-fibration 
\[
\varphi : V \to \bar{B}
\]
onto a smooth projective curve $\bar{B}$, where $B \subset \bar{B}$ is the smooth completion.

Since $X$ is $\Q$-factorial, the Picard group $\Pic(V)$ is finitely generated, hence $b_1(V)=0$. The general fiber of $\varphi$ is isomorphic to $\PP^1$, which is connected, and therefore the induced map $\varphi_\ast : \pi_1(V) \to \pi_1(\bar{B})$ and hence
\[
\varphi_\ast : H_1(V; \Z) \to H_1(\bar{B}; \Z)
\]
is surjective. Consequently, $b_1(\bar{B}) \le b_1(V) = 0$, implying $\bar{B} \cong \PP^1$. Thus $B$ is a smooth rational curve.

Since $\bar{B} \cong \PP^1$, one general fiber $F$ of $\varphi$, together with a cross-section, say $S$ of $\varphi$ and all but one irreducible component of every singular fiber, freely generate $\Pic(V)$.  

If $B$ is projective, then $\bar{B} \cong \PP^1$. Then a similar argument shows that a general fiber of $f$, together with all but one irreducible component of each singular fiber of $f$, freely generates $\Pic(X)$. Hence $\rho(X) \ge 1$, and equality holds if and only if $f$ has no reducible fibers. Therefore, $\Pic(X)$ cannot be torsion, contradicting the $\Q$-factoriality of $X$. This contradiction shows that $B$ must be affine, proving (1).

Now assume that $B$ is a smooth affine rational curve. By the same reasoning as above, all but one irreducible component of each singular fiber of $f$ freely generate $\Pic(X)$. Hence $\rho(X) \ge 1$ whenever $f$ has a reducible fiber, again contradicting the $\Q$-factoriality of $X$. Therefore, all fibers of $f$ are irreducible, establishing (2).
\end{proof}

\end{document}
